\documentclass[10pt,a4paper]{article}
\usepackage[margin=2.2cm]{geometry}
\usepackage{amsmath,amssymb,amsthm,booktabs,graphicx,microtype,siunitx}
\usepackage[table]{xcolor}
\usepackage[round,authoryear]{natbib}
\usepackage[hidelinks]{hyperref}
\usepackage{caption}
\usepackage{times}
\usepackage[section]{placeins}
\captionsetup{font=small,labelfont=bf}
\sisetup{detect-all}
\newtheorem{proposition}{Proposition}
\newcommand{\zv}{\mathbf{z}}
\input{macros.tex}
\input{macros_mb.tex}

\title{\vspace{-1.3cm}\textbf{Where Does the Gain in Dynamic Convolution Come From?}\\[3pt]
\large A controlled study of routing, capacity and width on two backbones\thanks{Draft of 17 October 2026. ResNet-18 and MobileNetV2 are complete; the same program on VGG-16, ResNet-50 and ConvNeXt-Tiny is pre-registered and running.}}
\author{Ankit Kumar\\\small Mechatronics Department, University of Agder, Grimstad, Norway}
\date{\small 24 September 2026; revised 17 October 2026}

\begin{document}
\maketitle
\vspace{-0.6cm}

\begin{abstract}
\noindent
Dynamic convolution replaces one kernel with a bank of $K$ kernels mixed
by input-dependent weights. Its gains are usually credited to
input-adaptive kernel selection, but published comparisons set it
against a network with $K$ times fewer parameters, so routing and
capacity are confounded. We separate them with a constant-mixture
control, a mismatched-input control and a frozen random router, on
ResNet-18 and MobileNetV2, CIFAR-100 and Tiny-ImageNet, across
\V{totruns} training runs. Two results hold on both backbones. The kernel
bank on its own adds nothing measurable: a uniformly mixed bank is never
significantly different from a plain network, and is statistically
equivalent to one at a $\pm$\SI{0.5}{pp} margin in \V{bothbankeq} of ten
conditions; the other three, all on MobileNetV2, are inconclusive. And
where dynamic convolution helps, routing carries the effect, while no
bank size beyond $K=2$ is distinguishable from $K=2$. Three ResNet-18
results do not transfer. On ResNet-18 a plain network of twice the width
beats the routed one at matched parameters in all three pairs, by 2.3 to
\SI{4.3}{pp}; on MobileNetV2 the routed network is not distinguishable
from a plain network with about twice its parameters in two of three
pairs. On ResNet-18 the routing effect vanishes on CIFAR-100 at width 64;
on MobileNetV2 it is \V{mbeffCC}\,pp there. And on MobileNetV2, annealing
the softmax temperature, which DY-Conv presents as necessary, costs about
\SI{1}{pp} at every width. MobileNetV2 adds a result with no ResNet-18
counterpart: a router frozen at its random initialisation is not
distinguishable from a trained one and delivers the full benefit over the
constant mixture, while routing each image with another image's
descriptor removes it. Routing entropy, measured on ResNet-18, does not
track the benefit. Of the four predictions we recorded before running the
experiments that test them, three failed, and we report each one.
\end{abstract}

\section{Introduction}

A dynamic convolution layer stores $K$ kernels and combines them per
input. For an input feature map $F$ and kernel bank
$\{\mathcal{W}_k\}_{k=1}^{K}$, the layer computes
\begin{equation}
\mathcal{W}_{\text{dyn}}(F)=\sum_{k=1}^{K} z_k(F)\,\mathcal{W}_k,
\qquad
\zv=\operatorname{softmax}\!\left(\frac{g(\phi(F))}{\tau}\right),
\label{eq:dynconv}
\end{equation}
where $\phi$ summarises the feature map, $g$ is a small multilayer
perceptron and $\tau$ is a softmax temperature. Nearly every published
variant sets $\phi$ to global average pooling
\citep{yang2019condconv,chen2020dynamic}.

The usual account of why this helps is that $\zv$ adapts the kernel to
the input. A second account is available and is rarely tested. The bank
multiplies the stored convolution parameters by $K$, and published
comparisons are made against a network with $K$ times fewer parameters.
Under that comparison, a gain from routing and a gain from extra
capacity look identical.

This report separates the two. We compare the standard routed layer
against a constant mixture that has the same bank and no routing, a
router fed a different image's descriptor, a router frozen at random
initialisation, and a plain network. The comparison spans three widths,
three bank sizes, two temperature schedules, two datasets and two
backbones: ResNet-18, where almost every parameter sits in a dynamic
layer, and MobileNetV2, where the dynamic layers are depthwise and hold
\V{mbdynshare}\,\% of the parameters. That is \V{totruns} training runs and
\V{totgpuh} GPU-hours on one RTX 4070 Ti. Every comparison is paired by
seed and tested for equivalence as well as for difference, so that a null
result is a claim with a stated margin rather than a failure to reject.
The MobileNetV2 program repeats the ResNet-18 program exactly, and its
predictions were recorded, dated, before its first run.

Two findings hold on both backbones:
\begin{enumerate}\setlength\itemsep{1pt}
\item The kernel bank adds nothing measurable on its own
(Sections~\ref{sec:decomp} and~\ref{sec:mbdecomp}).
\item Where dynamic convolution helps, the routing term carries the
effect, and banks larger than $K=2$ add nothing further
(Sections~\ref{sec:ksweep} and~\ref{sec:mbtemp}).
\end{enumerate}
Four depend on the backbone (Section~\ref{sec:transfer}):
\begin{enumerate}\setlength\itemsep{1pt}
\item Whether width beats routing at equal parameters. It does on
ResNet-18 in all three pairs we could form; on MobileNetV2 only at the
smallest width (Sections~\ref{sec:matched} and~\ref{sec:mbmatched}).
\item Where routing helps. On ResNet-18 not on CIFAR-100 at width 64; on
MobileNetV2 in all five conditions.
\item What temperature annealing does. On ResNet-18 it is neutral or
harmful; on MobileNetV2 it is harmful at every width.
\item Whether the router has to be trained. On MobileNetV2 a router
frozen at random initialisation does as well as a trained one
(Section~\ref{sec:mbw64}); on ResNet-18 the test is uninformative.
\end{enumerate}
On ResNet-18, routing entropy does not measure whether routing is useful
(Section~\ref{sec:entropy}).

We recorded four predictions before running the experiments that test
them. On ResNet-18 we attributed the CIFAR-100 width-64 null to the
network memorising its training set and predicted the effect would
vanish on Tiny-ImageNet at width 64; it did not (Section~\ref{sec:tin}).
For MobileNetV2 we predicted that the bank term would be equivalent to
zero in all five conditions, that width would beat routing in all three
pairs, and that routing would help in at least three conditions. Only the
last held (Section~\ref{sec:mbprereg}).

\section{Related work}

\paragraph{Dynamic convolution.} CondConv \citep{yang2019condconv} and
DY-Conv \citep{chen2020dynamic} introduced input-dependent mixtures of
kernels. ODConv \citep{li2022odconv} extends the attention across several
kernel dimensions, and KernelWarehouse \citep{li2024kernelwarehouse}
reorganises the bank so that far larger $K$ becomes affordable. The
reported gains are large: for a ResNet-18 on ImageNet at four times the
kernel parameters, \citet{li2024kernelwarehouse} report top-1
improvements of \SI{3.38}{pp} for DY-Conv, \SI{4.01}{pp} for ODConv and
\SI{5.61}{pp} for KernelWarehouse under their training recipe. FDConv
\citep{chen2025fdconv} observes that the parallel kernels of DY-Conv tend
toward highly similar frequency responses, so that the extra parameters
buy little adaptability. SD-Conv \citep{he2023sdconv} prunes half the
parameters of a trained dynamic network at a cost below \SI{1}{pp} on
CIFAR-100, and \citet{li2021dcd} argue that the $K$ static kernels are
redundant. To our knowledge none of these works isolates the bank from the
routing with a constant mixture of the same bank. Our finding that the
bank contributes nothing on its own is consistent with these
observations. Spatially resolved variants
such as DRConv \citep{chen2021drconv} route on local rather than global
signals and fall outside the controls used here.

\paragraph{Optimisation difficulty.} \citet{li2021dcd} identify the joint
optimisation of dynamic attention and static kernels as the central
difficulty of dynamic convolution and propose a matrix decomposition that
eases it. DY-Conv addresses the same difficulty by annealing the softmax
temperature from 30 to 1 early in training \citep{chen2020dynamic}.
Section~\ref{sec:temp} tests that schedule.

\paragraph{Re-parameterisation.} A constant mixture of $K$ kernels is a
linear re-parameterisation of a single kernel. Such re-parameterisations
can change optimisation without changing the function class
\citep{arora2018optimization}, and ACNet and RepVGG exploit this by
training a richer block and merging it into one kernel at inference
\citep{ding2019acnet,ding2021repvgg}. Our \texttt{uniform} arm is exactly
such a mergeable block.

\paragraph{Memorisation.} Networks of the size used here can fit their
training sets completely, even under random labels \citep{zhang2017}.
That motivated the saturation hypothesis tested in Section~\ref{sec:tin}.

\paragraph{Equivalence testing.} The two one-sided tests procedure
\citep{schuirmann1987} asserts that a difference lies within a stated
margin; \citet{lakens2017} give a practical treatment. It is standard in
bioequivalence studies and rare in this literature, where a large
$p$-value from a difference test is often read as evidence of no effect.

\section{Method}\label{sec:method}

\subsection{Routing signals}

All arms share Eq.~\eqref{eq:dynconv}, the backbone, the optimiser and the
seeds, and differ only in how $\zv$ is produced. Write
$\mu(F)\in\mathbb{R}^{C}$ for the channel-wise global average.

\begin{description}\setlength\itemsep{2pt}
\item[\texttt{static}] No dynamic layer: the plain backbone, $K=1$.
\item[\texttt{uniform}] $g\equiv\mathbf{0}$, so $\zv=(1/K,\dots,1/K)$ for
every input. The bank exists and is trained; nothing is routed.
\item[\texttt{gap}] $\zv=\operatorname{softmax}(g(\mu(F))/\tau)$ with $g$
trained. The standard layer.
\item[\texttt{shuffled}] As \texttt{gap}, but sample $i$ is routed with
the descriptor of sample $\pi(i)$ for a permutation $\pi$ redrawn on every
forward pass, at evaluation as well as training. The distribution of
$\zv$ is preserved and the pairing between image and kernel is
destroyed.
\item[\texttt{frozen\_random}] As \texttt{gap}, with every weight of $g$
frozen at its random initialisation. Routing varies with the input but
is never optimised.
\end{description}

\subsection{The constant mixture is a single kernel}

\begin{proposition}\label{prop:collapse}
If $\zv=(1/K,\dots,1/K)$ for every input, the layer computes
$F*\bar{\mathcal{W}}$ with $\bar{\mathcal{W}}=\tfrac1K\sum_k\mathcal{W}_k$,
and its effective parameter count equals that of one static kernel.
\end{proposition}
\begin{proof}
The mixture in Eq.~\eqref{eq:dynconv} is constant, so
$\mathcal{W}_{\text{dyn}}=\bar{\mathcal{W}}$ for every $F$.
\end{proof}

At inference \texttt{uniform} is therefore a plain ResNet-18. At width 64
it stores 44.18M parameters and has an effective count of 11.22M. It
spans the same function class as \texttt{static}, so any difference
between the two is an effect of optimisation. That makes it the right
control for the extra parameters of the bank.

\subsection{Two exact implementations}

Convolution is linear in its weights, so
\begin{equation}
F*\Big(\sum_k z_k\mathcal{W}_k\Big)=\sum_k z_k\,(F*\mathcal{W}_k).
\label{eq:lin}
\end{equation}
The left side builds a per-sample kernel and runs one grouped
convolution; at width 64 and batch 128 its aggregated weights alone take
about \SI{1.2}{\giga\byte} per stage-4 layer. The right side runs $K$
ordinary convolutions and mixes the outputs, which costs $K$ times the
arithmetic but no per-sample weights. Every run uses the right side. The
test suite asserts that the two agree to $10^{-10}$ in double precision
across four routing signals.

\subsection{Training protocol}

\paragraph{Data.} CIFAR-100 \citep{krizhevsky2009}, $32\times32$, 100
classes, at widths 16, 32 and 64. Tiny-ImageNet \citep{le2015tiny},
$64\times64$, 200 classes, at widths 32 and 64. The stem is the same on
both: a $3\times3$ stride-1 convolution with no max-pool. At
$64\times64$ this leaves an $8\times8$ final feature map rather than
$4\times4$ and roughly quadruples the cost of the deeper stages, but
changing the stem would change the network between datasets. Crop
padding is 4 pixels at 32 and 8 at 64.

\paragraph{Network and recipe.} ResNet-18 \citep{he2016deep} in which all
16 $3\times3$ convolutions inside residual blocks are dynamic; the stem
and the $1\times1$ shortcuts stay static. $K=4$ unless stated. SGD with
Nesterov momentum 0.9, learning rate 0.1, weight decay $5\times10^{-4}$,
batch 128, five warm-up epochs then cosine decay over 150 epochs.
Augmentation is random cropping and horizontal flipping only. Mixed
precision in bfloat16. Six seeds, 0 to 5, shared by every arm. Batch 128
held at Tiny-ImageNet width 64 after a one-epoch memory probe, so the two
Tiny-ImageNet widths are directly comparable.

\paragraph{Second backbone.} MobileNetV2 \citep{sandler2018mobilenetv2},
adapted to small inputs in the usual way: the stem and the second stage
use stride 1. Width 64 is the standard network; widths 32 and 16 scale
every channel count by $1/2$ and $1/4$, rounded to a multiple of 8. The
rule that decides which layers are dynamic is copied from ResNet-18
without change: every $k\times k$ convolution with $k>1$ inside a block.
In MobileNetV2 that is the 17 depthwise $3\times3$ convolutions; the
stem, every $1\times1$ expansion and projection, and the classifier stay
static. The consequences differ from ResNet-18. At CIFAR-100 width 64 the
bank adds \V{mbbankadd}M parameters to a \V{mbpstatic}M network, so only
\V{mbdynshare}\,\% of the static parameters sit in dynamic layers, and the
router of the \texttt{gap} arm (the ResNet-18 design, a two-layer
perceptron of hidden width $C_{\text{in}}/4$) has \V{mbrouterK}K
parameters, \V{mbrouterpct}\,\% of the static network and more than the
bank it routes. Static convolutions and every bank kernel share one
initialisation, Kaiming-normal with fan-out. Recipe, arms, seeds and
margin are those of ResNet-18; every group ran at batch 128.

\paragraph{Temperature.} $\tau=1$ throughout, except in
Section~\ref{sec:temp}, where $\tau$ falls linearly from 30 to 1 over the
first ten epochs and stays at 1 thereafter.

\section{Statistical protocol}\label{sec:stats}

\paragraph{Equivalence.} Several conclusions are null conclusions, and a
large $p$ from a difference test does not establish one. For paired
differences $d_i$ with mean $\bar d$ and standard error $s_{\bar d}$ we
use two one-sided tests against a margin $\delta$,
\begin{equation}
p_{\text{TOST}}=\max\Big\{1-T_{n-1}\!\Big(\tfrac{\bar d+\delta}{s_{\bar d}}\Big),\;
T_{n-1}\!\Big(\tfrac{\bar d-\delta}{s_{\bar d}}\Big)\Big\},
\end{equation}
and read $p_{\text{TOST}}<0.05$ as a positive claim that the true
difference lies within $\pm\delta$. We fixed $\delta=\SI{0.5}{pp}$ before
the first run. The pre-registration justified this by reference to
reported gains at CIFAR scale without citing a source. The nearest
documented figures we have since found are the ImageNet ResNet-18 gains
of 3.38 to \SI{5.61}{pp} quoted above, and the margin sits an
order of magnitude below them.

\paragraph{Why six seeds.} The two-sided Wilcoxon signed-rank test
\citep{wilcoxon1945} on $n$ pairs admits only $2^n$ equally likely sign
assignments, so its smallest attainable $p$ is $2^{1-n}$: 0.25 at $n=3$,
0.0625 at $n=5$, 0.031 at $n=6$. Six is the smallest count at which it
can reach 0.05. We report it beside the paired $t$-test throughout.

\paragraph{Power.} For $\texttt{gap}-\texttt{uniform}$ on CIFAR-100 at
width 64 the paired standard deviation is \V{powsd}\,pp, giving a
minimum detectable difference of \V{mde}\,pp at $n=6$. Simulated
power ($2\times10^5$ draws) is \V{pwA}\,\% for a true effect of
\SI{0.25}{pp}, \V{pwB}\,\% at \SI{0.5}{pp} and
\V{pwC}\,\% at \SI{1.0}{pp}. The width-64 null on CIFAR-100
is not a power artefact.

\paragraph{Pre-registration.} Arms, recipe, primary metric, seeds, margin
and the decision table were written down before the first run finished.
The primary metric is final-epoch top-1 accuracy, since best-epoch
accuracy maximises over a noisy sequence. The Tiny-ImageNet width-64
prediction of Section~\ref{sec:tin} was recorded in writing, dated,
before any of those runs. An expectation about Tiny-ImageNet width 32 was
stated before those runs but not recorded; we give it correspondingly
less weight. For MobileNetV2, three predictions and a gate (the first
static run had to reach \SI{65}{\percent} on CIFAR-100 before the other
131 runs were spent) were recorded on 30 September 2026, before any
MobileNetV2 run (Section~\ref{sec:mbprereg}).

\section{ResNet-18 on CIFAR-100 at the standard width}\label{sec:w64}

Sections~\ref{sec:w64} to~\ref{sec:entropy} report ResNet-18.
Section~\ref{sec:mb} reports MobileNetV2 and Section~\ref{sec:transfer}
compares the two.

\begin{figure}[!htbp]\centering
\includegraphics[width=\textwidth]{fig/f3_w64.pdf}
\caption{CIFAR-100, width 64. (a) Every run of the five arms; bars are
arm means and the shaded band is $\pm$\SI{0.5}{pp} around
\texttt{static}. (b) Router flatness $1-\tilde H$ per seed, log scale;
\texttt{uniform} is exactly zero and is drawn at the axis floor.}
\label{fig:w64}
\end{figure}

At width 64 every arm is equivalent to a plain ResNet-18
(Table~\ref{tab:vsstatic}, Figure~\ref{fig:w64}). The largest gap between
arm means is \V{spread64}\,pp, against a median within-arm seed
standard deviation of \V{seedsd64}\,pp. A one-way ANOVA across the
five arms gives $F(4,25)=\V{anovaF}$, $p=\V{anovaP}$.

A router whose weights were never trained matches a trained one
($\texttt{frozen\_random}-\texttt{gap}=\V{frzVgap}$\,pp, $p=\V{frzVgapP}$,
$p_{\text{TOST}}=\V{frzVgapT}$), and so does a router that routes each
image with another image's descriptor. The trained router is not
degenerate: its per-seed mean entropy stays below that of the controls
(Figure~\ref{fig:w64}b). It discriminates between kernels, and at this
width the discrimination changes nothing.

\section{ResNet-18: the capacity sweep on CIFAR-100}\label{sec:capacity}

\begin{figure}[!htbp]\centering
\includegraphics[width=\textwidth]{fig/f1_sweep.pdf}
\caption{CIFAR-100. (a) Accuracy against base width, $\pm1$ standard
deviation over six seeds. (b) Paired difference from \texttt{static} with
\SI{95}{\percent} confidence intervals.}
\label{fig:sweep}
\end{figure}

Reducing width changes the answer (Figure~\ref{fig:sweep}).
$\texttt{gap}-\texttt{static}$ is \V{effCA}\,pp at width 16
($p\V{effCAp}$), \V{effCB}\,pp at width 32 ($p\V{effCBp}$) and
\V{effCC}\,pp at width 64 ($p\V{effCCp}$). The effect is not monotone in
capacity: it peaks at width 32. MobileNetV2 shows no such peak and no
null at width 64 (Section~\ref{sec:mbdecomp}). \texttt{uniform} is equivalent to
\texttt{static} at all three widths. Figure~\ref{fig:allruns} in the
appendix shows every run.

\section{ResNet-18: the kernel bank contributes nothing}\label{sec:decomp}

\begin{figure}[!htbp]\centering
\includegraphics[width=0.84\textwidth]{fig/f4_decomp.pdf}
\caption{Decomposition of the effect of dynamic convolution into the
bank ($\texttt{uniform}-\texttt{static}$) and routing
($\texttt{gap}-\texttt{uniform}$), five conditions, \SI{95}{\percent}
CIs.}
\label{fig:decomp}
\end{figure}

Splitting the total effect $\texttt{gap}-\texttt{static}$ into a bank
term and a routing term separates the two accounts of the introduction.
Table~\ref{tab:decomp} and Figure~\ref{fig:decomp} give the split in all
five conditions that have a \texttt{uniform} arm.

The bank term is equivalent to zero in every one of them, on both
datasets, with $p_{\text{TOST}}$ between 0.014 and 0.025. The largest
bank term observed is \V{bankCB}\,pp. Since the bank contributes nothing
within the margin, the total effect is the routing effect: on
Tiny-ImageNet at width 64, for example, routing is \V{routTB}\,pp of a
\V{effTB}\,pp total. At CIFAR-100 width 16 the bank term is slightly
negative (\V{bankCA}\,pp, equivalent) and the routing term therefore
exceeds the total, which is why we report no percentage share.

\begin{table}[!htbp]\centering\small
\caption{Bank and routing terms, paired by seed. The bank term is
equivalent to zero in all five conditions.}
\label{tab:decomp}
\begin{tabular}{l r rll rl rl}
\toprule
& & \multicolumn{3}{c}{bank (\texttt{uniform}$-$\texttt{static})}
& \multicolumn{2}{c}{routing (\texttt{gap}$-$\texttt{uniform})}
& \multicolumn{2}{c}{total} \\
\cmidrule(lr){3-5}\cmidrule(lr){6-7}\cmidrule(lr){8-9}
dataset & width & $\Delta$ (pp) & $p_{\text{TOST}}$ & verdict & $\Delta$ (pp) & $p$ & $\Delta$ (pp) & $p$ \\
\midrule
%%DECOMP%%
\bottomrule
\end{tabular}
\end{table}

\section{ResNet-18: bank size}\label{sec:ksweep}

\begin{figure}[!htbp]\centering
\includegraphics[width=0.84\textwidth]{fig/f6_k.pdf}
\caption{CIFAR-100, width 32. (a) Effect over \texttt{static} for
$K\in\{2,4,8\}$, \SI{95}{\percent} CIs. (b) Stored parameters and
training time.}
\label{fig:k}
\end{figure}

The width-32 peak was measured at $K=4$; Table~\ref{tab:k} repeats it at
$K=2$ and $K=8$. All three bank sizes beat \texttt{static} by a similar
margin, and no pair of them is distinguishable (Table~\ref{tab:kpairs}):
$K=8$ over $K=2$ is \V{k82}\,pp ($p=\V{k82p}$). The benefit is already
complete at $K=2$. Going from $K=2$ to $K=8$ multiplies stored parameters
by \V{kpar} and training time by \V{ktime} (Figure~\ref{fig:k}).

\begin{table}[!htbp]\centering\small
\caption{Bank size at CIFAR-100 width 32. Differences are against the
\texttt{static} row.}
\label{tab:k}
\begin{tabular}{l rr rl l rr}
\toprule
$K$ & accuracy & sd & $\Delta$ (pp) & \SI{95}{\percent} CI & $p$ & params (M) & min/run \\
\midrule
%%KSWEEP%%
\bottomrule
\end{tabular}
\end{table}

\begin{table}[!htbp]\centering\small
\caption{Pairwise comparisons between bank sizes.}
\label{tab:kpairs}
\begin{tabular}{l rl ll l}
\toprule
comparison & $\Delta$ (pp) & \SI{95}{\percent} CI & $p$ & $p_{\text{TOST}}$ & verdict \\
\midrule
%%KPAIRS%%
\bottomrule
\end{tabular}
\end{table}

\section{ResNet-18 on Tiny-ImageNet, and a failed prediction}\label{sec:tin}

\begin{figure}[!htbp]\centering
\includegraphics[width=\textwidth]{fig/f7_tin.pdf}
\caption{(a) All 36 Tiny-ImageNet runs. (b) $\texttt{gap}-\texttt{static}$
by width on both datasets, \SI{95}{\percent} CIs. The CIFAR-100
width-64 null does not appear on Tiny-ImageNet.}
\label{fig:tin}
\end{figure}

\subsection{Width 32}

On Tiny-ImageNet at width 32, $\texttt{gap}-\texttt{static}$ is
\V{effTA}\,pp (\SI{95}{\percent} CI [\V{effTAlo}, \V{effTAhi}],
$p\V{effTAp}$), against \V{effCB}\,pp on CIFAR-100 at the same width. The
bank term is \V{bankTA}\,pp and equivalent to zero. Before these runs we
expected the effect to be larger than on CIFAR-100, since the static arm
trains to a higher loss (\V{lossTA} against \V{lossCB}) and so sits
further from memorisation. That expectation held, but it was stated and
not recorded, it predicted only a direction, and we do not treat it as
strong evidence.

\subsection{Width 64: the recorded prediction}

The saturation hypothesis says the routing effect vanishes once a network
memorises its training data. On CIFAR-100 the width-64 null coincides
with a final training loss of \V{lossCC}. Before running Tiny-ImageNet at
width 64 we recorded, in writing and dated, the prediction in
Table~\ref{tab:predict}: that the effect would fall below \SI{0.5}{pp}
and be equivalent to zero, and that the static arm would train to a loss
below 0.05.

\begin{table}[!htbp]\centering\small
\caption{The prediction recorded before the Tiny-ImageNet width-64 runs,
and the outcome. Both conditions were stated as required.}
\label{tab:predict}
\begin{tabular}{l p{5.4cm} p{4.6cm} l}
\toprule
& prediction & measured & outcome \\
\midrule
%%PREDICT%%
\bottomrule
\end{tabular}
\end{table}

The loss condition held (\V{lossTB}) and the effect did not vanish:
$\texttt{gap}-\texttt{static}=\V{effTB}$\,pp, \SI{95}{\percent} CI
[\V{effTBlo}, \V{effTBhi}], $p=\V{predP}$. The prediction failed. The
bank term is again equivalent to zero (\V{bankTB}\,pp,
$p_{\text{TOST}}=\V{bankTBT}$), so the effect is carried by routing, and
it is as large at width 64 as at width 32.

\subsection{What remains of the saturation account}

\begin{figure}[!htbp]\centering
\includegraphics[width=\textwidth]{fig/f8_prediction.pdf}
\caption{(a) Training loss of static and gap arms, seed 0, log scale.
(b) Routing effect against the final training loss of the static arm in
all five conditions. The shaded band is the region the prediction
associated with a null.}
\label{fig:pred}
\end{figure}

Table~\ref{tab:satur} orders the five conditions by the final training
loss of the static arm. The effects do not follow that order
(Figure~\ref{fig:pred}b). Tiny-ImageNet at width 64 trains to a loss
close to CIFAR-100 at width 32 and shows twice the effect; CIFAR-100 at
width 16 trains to the highest loss and shows one of the smallest. A
threshold near 0.01 would separate the one null from the four positive
results, but it would be chosen after seeing the data, and we do not
adopt it.

What survives is narrower. On CIFAR-100 the routing effect vanishes at
width 64, and that result is well powered and holds under two temperature
schedules. It does not replicate on Tiny-ImageNet, and training loss does
not explain it. Plausible differences between the datasets include the
input resolution, the final feature-map size and the number of classes;
none of them was isolated here.

\begin{table}[!htbp]\centering\small
\caption{Five conditions ordered by the final training loss of the
static arm.}
\label{tab:satur}
\begin{tabular}{l r r rl l}
\toprule
dataset & width & static train loss & $\texttt{gap}-\texttt{static}$ (pp) & \SI{95}{\percent} CI & verdict \\
\midrule
%%SATUR%%
\bottomrule
\end{tabular}
\end{table}

\section{ResNet-18: matched parameter budgets}\label{sec:matched}

\begin{figure}[!htbp]\centering
\includegraphics[width=\textwidth]{fig/f5_matched.pdf}
\caption{Accuracy against stored parameters. Dotted arrows join routed
and plain networks matched to within \SI{1}{\percent} on parameter
count.}
\label{fig:matched}
\end{figure}

Because widths differ by a factor of two and the bank multiplies
convolution parameters by four, a routed network at one width stores
almost exactly as many parameters as a plain network at the next width
up. That gives three matched comparisons without any extra runs
(Table~\ref{tab:matched}, Figure~\ref{fig:matched}).

The plain network wins all three: by \V{mMA}\,pp at 2.8M
parameters and \V{mMB}\,pp at 11.2M on CIFAR-100, and by
\V{mMC}\,pp at 11.2M on Tiny-ImageNet. It also trains faster in
each case; on Tiny-ImageNet, \V{tMCs}\,min per run against
\V{tMCg}\,min. A fourth, unmatched comparison points the same
way: the $K=8$ network at width 32 stores 22.1M parameters, twice the
plain width-64 network, and trails it by \V{k8vs64}\,pp. This result
depends on almost every parameter of ResNet-18 sitting in a dynamic
layer, and it does not transfer to MobileNetV2
(Section~\ref{sec:mbmatched}).

\begin{table}[!htbp]\centering\small
\caption{Routed against plain networks at matched parameter counts,
paired by seed. Minutes are routed / plain per run. The last column is
the routed network's stored parameters divided by the plain one's; a pair
is called matched only within 10\,\%.}
\label{tab:matched}
\resizebox{\textwidth}{!}{%
\begin{tabular}{l l l rr rl l r l}
\toprule
dataset & routed & plain & routed & plain & $\Delta$ (pp) & \SI{95}{\percent} CI & $p$ & min/run & param ratio \\
\midrule
%%MATCHED%%
\bottomrule
\end{tabular}}
\end{table}

\section{ResNet-18: temperature schedule}\label{sec:temp}

\begin{figure}[!htbp]\centering
\includegraphics[width=\textwidth]{fig/f9_temp.pdf}
\caption{CIFAR-100. (a) Routing effect under a constant and an annealed
temperature, \SI{95}{\percent} CIs. (b) Router flatness. (c) Flatness
against benefit over \texttt{uniform}; the two schedules do not lie on
one curve.}
\label{fig:temp}
\end{figure}

With $\tau=1$ the softmax can be close to one-hot from the first step,
starving most of the bank of gradient. DY-Conv anneals $\tau$ from 30 to
1 for this reason \citep{chen2020dynamic}. If the CIFAR-100 width-64
null came from an undertrained bank, it would say nothing about dynamic
convolution, so we repeated \texttt{gap} at all three widths with the
annealed schedule (Table~\ref{tab:temp}).

The width-64 null survives: annealed \texttt{gap} against \texttt{static}
is \V{ansC}\,pp with $p_{\text{TOST}}=\V{ansCT}$, a tighter equivalence
than the \V{constT64} obtained at $\tau=1$. Annealing is otherwise
neutral at width 64 (\V{anC}\,pp, $p=\V{anCp}$) and width 16
(\V{anA}\,pp, $p=\V{anAp}$), and harmful at width 32 (\V{anB}\,pp,
$p=\V{anBp}$), where it also leaves a higher final training loss
(\V{laB} against \V{lcB}). Backbone, dataset, bank size and schedule
length all differ from the DY-Conv setting, so this is not evidence
against their result; it is evidence that the CIFAR-100 null does not
depend on the choice of schedule. On MobileNetV2 annealing is harmful at
all three widths (Section~\ref{sec:mbtemp}).

At width 64 the annealed router went flat during the high-$\tau$ phase
and stayed flat for the remaining 140 epochs: it is \V{flC} times flatter
than the constant router ($1-\tilde H=\V{haC}$) and equivalent to
\texttt{uniform} (\V{anuC}\,pp, $p_{\text{TOST}}=\V{anuCT}$). This follows
from Proposition~\ref{prop:collapse}. A flat mixture trains the bank
toward its own mean, and a homogeneous bank leaves the router nothing to
discriminate between. At width 16 the same schedule did not produce that
outcome, which the next section takes up.

\begin{table}[!htbp]\centering\small
\caption{Temperature ablation on CIFAR-100, \texttt{gap} arm. The last
two columns test whether the annealed router behaves like the constant
mixture.}
\label{tab:temp}
\begin{tabular}{r rr rl r rl rl}
\toprule
& \multicolumn{2}{c}{accuracy} & \multicolumn{2}{c}{annealed $-$ constant}
& \multicolumn{3}{c}{vs \texttt{static}} & \multicolumn{2}{c}{annealed vs \texttt{uniform}} \\
\cmidrule(lr){2-3}\cmidrule(lr){4-5}\cmidrule(lr){6-8}\cmidrule(lr){9-10}
width & $\tau=1$ & annealed & $\Delta$ & $p$ & $\tau=1$ & annealed & $p_{\text{TOST}}$ & $\Delta$ & verdict \\
\midrule
%%TEMP%%
\bottomrule
\end{tabular}
\end{table}

\section{ResNet-18: routing entropy does not measure routing usefulness}\label{sec:entropy}

\begin{figure}[!htbp]\centering
\includegraphics[width=0.74\textwidth]{fig/f10_entropy.pdf}
\caption{CIFAR-100, width 16. The annealed router is \V{flA} times flatter
by entropy and delivers the larger benefit over \texttt{uniform}.}
\label{fig:ent}
\end{figure}

Normalised routing entropy $\tilde H=H(\zv)/\log K$ is the statistic
this literature reports to show that a router is doing something.
Table~\ref{tab:entben} sets it against the benefit each router delivers
over \texttt{uniform} in six conditions.

At width 16 the two schedules give a counterexample
(Figure~\ref{fig:ent}). The constant router has $1-\tilde H=\V{hcA}$ and
beats \texttt{uniform} by \V{bcA}\,pp. The annealed router has
$1-\tilde H=\V{haA}$, \V{flA} times flatter, and beats \texttt{uniform} by
\V{anuA}\,pp ($p=\V{anuAp}$). A router that looks nearly inactive on the
standard diagnostic delivered more than one that looks clearly active.

The reason is that entropy compresses the signal quadratically. Writing
$z_k=1/K+\varepsilon_k$ with $\sum_k\varepsilon_k=0$ and expanding to
second order,
\begin{equation}
1-\tilde H\;\approx\;\frac{K}{2\log K}\sum_k\varepsilon_k^{2},
\label{eq:ent}
\end{equation}
so the deviations that carry input-dependent information enter squared.
At $1-\tilde H=\V{haA}$ the root-mean-square deviation from $1/K$ is
about 0.022, so typical mixing weights lie in
[\V{spanlo}, \V{spanhi}], roughly \V{spanpct}\,\% of $1/K$.
Entropy reports that as 0.997.

Entropy is not meaningless. The annealed router at width 64, at
$1-\tilde H=\V{haC}$, has collapsed in the sense that matters: it is
equivalent to \texttt{uniform}. But across the six conditions of
Table~\ref{tab:entben}, sorting by $1-\tilde H$ does not sort the
benefit, and the two schedules do not lie on one curve
(Figure~\ref{fig:temp}c). A value of $\tilde H$ below one is not evidence
that routing is useful, and a value near one is not evidence that it is
not.

\begin{table}[!htbp]\centering\small
\caption{Routing entropy against the benefit it is taken to indicate,
CIFAR-100. RMS dev is the root-mean-square deviation of the mixing
weights from $1/K$, recovered from Eq.~\eqref{eq:ent}.}
\label{tab:entben}
\begin{tabular}{r l rrr rl}
\toprule
width & $\tau$ & $\tilde H$ & $1-\tilde H$ & RMS dev & \texttt{gap}$-$\texttt{uniform} (pp) & $p$ \\
\midrule
%%ENTBEN%%
\bottomrule
\end{tabular}
\end{table}

\section{A second backbone: MobileNetV2}\label{sec:mb}

\subsection{Setup and coverage}

The program of Sections~\ref{sec:w64} to~\ref{sec:entropy} was repeated
on MobileNetV2 with no change to arms, seeds, recipe, margin or the rule
that decides which layers are dynamic (Section~\ref{sec:method}). The
gate run reached \V{mbgate}\,\% against the recorded floor of
\SI{65}{\percent}, and all \V{mbruns} runs finished in \V{mbgpuh}
GPU-hours. None diverged and every group trained at batch 128. One run
(Tiny-ImageNet, width 64, \texttt{gap}, seed 1, final accuracy
\V{mbresumedacc}\,\%) was stopped by a GPU fault at epoch 125 and resumed
from its checkpoint; without it, $\texttt{gap}-\texttt{static}$ in that
condition is \V{mbsensTB}\,pp ($p=\V{mbsensTBp}$, $n=5$), so no conclusion
depends on it. The MobileNetV2 numbers in this draft are taken from the
training log, which prints every run's final accuracy and training loss;
the routing entropies are recorded only in the result files and are not
analysed here.

\subsection{The bank adds nothing measurable; routing helps everywhere}\label{sec:mbdecomp}

\begin{figure}[!htbp]\centering
\includegraphics[width=\textwidth]{fig/g1_decomp_both.pdf}
\caption{Bank term ($\texttt{uniform}-\texttt{static}$) and routing term
($\texttt{gap}-\texttt{uniform}$) in the five conditions, for (a)
ResNet-18 and (b) MobileNetV2, \SI{95}{\percent} CIs. The shaded band is
the $\pm$\SI{0.5}{pp} margin.}
\label{fig:mbdecomp}
\end{figure}

On MobileNetV2 the bank term is never significantly different from zero
(Table~\ref{tab:mbdecomp}, Figure~\ref{fig:mbdecomp}). It is equivalent
to zero on CIFAR-100 at widths 32 and 64 ($p_{\text{TOST}}=\V{mbbankCBT}$
and \V{mbbankCCT}). At CIFAR-100 width 16 and at both Tiny-ImageNet widths
it is inconclusive: the point estimates (\V{mbbankCA}, \V{mbbankTA} and
\V{mbbankTB}\,pp) lie inside the margin, but the confidence intervals do
not. At CIFAR-100 width 16 the static arm itself is unusually variable
across seeds (Table~\ref{tab:mballruns}).

The routing term is positive with $p<0.05$ in all five conditions, from
\V{mbroutmin} to \V{mbroutmax}\,pp, and the total effect ranges from
\V{mbtotmin} to \V{mbtotmax}\,pp. Unlike ResNet-18, the effect does not
vanish on CIFAR-100 at width 64: $\texttt{gap}-\texttt{static}$ is
\V{mbeffCC}\,pp there (\SI{95}{\percent} CI [\V{mbeffCClo},
\V{mbeffCChi}], $p=\V{mbeffCCp}$), and the effect is roughly flat across
width on both datasets (Figure~\ref{fig:width}).

\begin{table}[!htbp]\centering\small
\caption{MobileNetV2: bank and routing terms, paired by seed.}
\label{tab:mbdecomp}
\begin{tabular}{l r rll rl rl}
\toprule
& & \multicolumn{3}{c}{bank (\texttt{uniform}$-$\texttt{static})}
& \multicolumn{2}{c}{routing (\texttt{gap}$-$\texttt{uniform})}
& \multicolumn{2}{c}{total} \\
\cmidrule(lr){3-5}\cmidrule(lr){6-7}\cmidrule(lr){8-9}
dataset & width & $\Delta$ (pp) & $p_{\text{TOST}}$ & verdict & $\Delta$ (pp) & $p$ & $\Delta$ (pp) & $p$ \\
\midrule
%%MBDECOMP%%
\bottomrule
\end{tabular}
\end{table}

\begin{figure}[!htbp]\centering
\includegraphics[width=\textwidth]{fig/g3_width_both.pdf}
\caption{$\texttt{gap}-\texttt{static}$ against width for both backbones,
\SI{95}{\percent} CIs. ResNet-18 has a null on CIFAR-100 at width 64;
MobileNetV2 does not.}
\label{fig:width}
\end{figure}

\subsection{An untrained router does as well as a trained one}\label{sec:mbw64}

\begin{figure}[!htbp]\centering
\includegraphics[width=\textwidth]{fig/g2_mb_w64.pdf}
\caption{CIFAR-100, width 64. (a) Every MobileNetV2 run of the five arms;
bars are arm means, the band is $\pm$\SI{0.5}{pp} around \texttt{static}.
(b) The routing controls on both backbones, \SI{95}{\percent} CIs.}
\label{fig:mbw64}
\end{figure}

On CIFAR-100 at width 64 the five MobileNetV2 arms differ (one-way ANOVA
$F=\V{mbanovaF}$, $p\V{mbanovaP}$), where the ResNet-18 arms did not
($F=\V{anovaF}$, $p=\V{anovaP}$). The controls separate cleanly
(Table~\ref{tab:mbw64}, Figure~\ref{fig:mbw64}).

A router frozen at its random initialisation beats \texttt{static} by
\V{mbwfrozenrandomstatic}\,pp and \texttt{uniform} by
\V{mbwfrozenrandomuniform}\,pp ($p=\V{mbwfrozenrandomuniformp}$), and is
not distinguishable from the trained router
($\texttt{frozen\_random}-\texttt{gap}=\V{mbwfrozenrandomgap}$\,pp,
\SI{95}{\percent} CI [\V{mbwfrozenrandomgaplo},
\V{mbwfrozenrandomgaphi}], $p=\V{mbwfrozenrandomgapp}$). It is not
equivalent to it at our margin either ($p_{\text{TOST}}=\V{mbwfrozenrandomgapT}$),
so the claim is that no difference was detected, not that none exists.
Routing each image with another image's descriptor removes the benefit:
\texttt{shuffled} is not distinguishable from \texttt{static}
(\V{mbwshuffledstatic}\,pp, $p=\V{mbwshuffledstaticp}$) and falls below
\texttt{gap} by \V{mbwshuffledgap}\,pp ($p=\V{mbwshuffledgapp}$).

Read together, the benefit on MobileNetV2 requires that each image's own
descriptor determine its mixture, and does not require the router to be
trained. A fixed random map from descriptors to mixtures is enough; the
bank learns kernels that suit the routing it is given. The frozen
router's weights are held as buffers and are not trained, so the gain is
not a product of the \V{mbrouterK}K trainable router parameters either.
On ResNet-18 the same controls are uninformative, since no arm differed
from \texttt{static} at this width. The frozen router was run only on
CIFAR-100 at width 64, and the claim is limited to that condition.

\begin{table}[!htbp]\centering\small
\caption{Routing controls on CIFAR-100 at width 64. The ResNet-18 columns
are repeated from Section~\ref{sec:w64} for comparison.}
\label{tab:mbw64}
\resizebox{\textwidth}{!}{%
\begin{tabular}{l rl rl l l l}
\toprule
& \multicolumn{2}{c}{ResNet-18} & \multicolumn{5}{c}{MobileNetV2} \\
\cmidrule(lr){2-3}\cmidrule(lr){4-8}
comparison & $\Delta$ (pp) & verdict & $\Delta$ (pp) & \SI{95}{\percent} CI & $p$ & $p_{\text{TOST}}$ & verdict \\
\midrule
%%MBW64%%
\bottomrule
\end{tabular}}
\end{table}

\subsection{Parameter budgets}\label{sec:mbmatched}

\begin{figure}[!htbp]\centering
\includegraphics[width=\textwidth]{fig/g4_mb_matched.pdf}
\caption{MobileNetV2: accuracy against stored parameters. Arrows join
\texttt{gap} at width $c$ to \texttt{static} at width $2c$; labels give
the accuracy difference and the parameter ratio.}
\label{fig:mbmatched}
\end{figure}

On ResNet-18, \texttt{gap} at width $c$ and \texttt{static} at width $2c$
store nearly the same number of parameters. On MobileNetV2 they do not:
because so few parameters sit in dynamic layers, the routed network has
\V{mbmMAr} to \V{mbmMBr} times the parameters of the plain one
(Table~\ref{tab:mbmatched}, Figure~\ref{fig:mbmatched}). Even so, the
plain network with roughly twice the parameters wins only at the smallest
width (\V{mbmMA}\,pp, $p=\V{mbmMAp}$). On CIFAR-100 at width 32 and on
Tiny-ImageNet the two are not distinguishable (\V{mbmMB}\,pp,
$p=\V{mbmMBp}$; \V{mbmMC}\,pp, $p=\V{mbmMCp}$). Interpolating the static
accuracy curve in log-parameters to the routed network's parameter count
puts the routed network ahead in all three cases (static minus routed:
\V{mbmMAi}, \V{mbmMBi} and \V{mbmMCi}\,pp); that comparison is an
interpolation and has no confidence interval.

\begin{table}[!htbp]\centering\small
\caption{MobileNetV2: \texttt{gap} at width $c$ against \texttt{static} at
width $2c$, paired by seed. $\Delta$ is plain minus routed. The ratio is
the routed network's stored parameters over the plain one's; none of
these pairs is matched. The last column compares \texttt{gap} with the
static curve interpolated to its parameter count.}
\label{tab:mbmatched}
\resizebox{\textwidth}{!}{%
\begin{tabular}{l l l rr rl l l r r}
\toprule
dataset & routed & plain & routed & plain & $\Delta$ (pp) & \SI{95}{\percent} CI & $p$ & verdict & ratio & $\Delta$ equal params \\
\midrule
%%MBMATCHED%%
\bottomrule
\end{tabular}}
\end{table}

\subsection{Bank size and temperature}\label{sec:mbtemp}

\begin{figure}[!htbp]\centering
\includegraphics[width=\textwidth]{fig/g5_temp_k.pdf}
\caption{Both backbones, CIFAR-100. (a) Annealed minus constant
temperature for \texttt{gap}. (b) $\texttt{gap}-\texttt{static}$ at width 32
for $K\in\{2,4,8\}$. \SI{95}{\percent} CIs.}
\label{fig:mbtemp}
\end{figure}

Bank size behaves as on ResNet-18 (Table~\ref{tab:mbk},
Figure~\ref{fig:mbtemp}b): all three sizes beat \texttt{static}, and
$K=8$ over $K=2$ is \V{mbk82}\,pp ($p=\V{mbk82p}$, \V{mbk82v}).

Temperature does not (Table~\ref{tab:mbtemp}, Figure~\ref{fig:mbtemp}a).
Annealing $\tau$ from 30 to 1 lowers \texttt{gap} at every width:
\V{mbanA}\,pp at width 16 ($p=\V{mbanAp}$), \V{mbanB}\,pp at width 32
($p=\V{mbanBp}$) and \V{mbanC}\,pp at width 64 ($p=\V{mbanCp}$). At width
32 the annealed router is equivalent to the constant mixture
(\V{mbanuB}\,pp); at width 64 it remains above it (\V{mbanuC}\,pp). The
annealed runs also end at a higher training loss at every width
(\V{mblaA}, \V{mblaB}, \V{mblaC} against \V{mblcA}, \V{mblcB},
\V{mblcC}). On ResNet-18 at width 64 the annealed router went flat and
stayed flat (Section~\ref{sec:temp}); whether the same happens on
MobileNetV2 needs the routing entropies, which this draft does not
include.

\begin{table}[!htbp]\centering\small
\caption{MobileNetV2: bank size at CIFAR-100 width 32.}
\label{tab:mbk}
\begin{tabular}{l rr rl l rr}
\toprule
$K$ & accuracy & sd & $\Delta$ (pp) & \SI{95}{\percent} CI & $p$ & params (M) & min/run \\
\midrule
%%MBK%%
\bottomrule
\end{tabular}
\end{table}

\begin{table}[!htbp]\centering\small
\caption{MobileNetV2: temperature ablation on CIFAR-100, \texttt{gap} arm.}
\label{tab:mbtemp}
\begin{tabular}{r rr rl r rl rl}
\toprule
& \multicolumn{2}{c}{accuracy} & \multicolumn{2}{c}{annealed $-$ constant}
& \multicolumn{3}{c}{vs \texttt{static}} & \multicolumn{2}{c}{annealed vs \texttt{uniform}} \\
\cmidrule(lr){2-3}\cmidrule(lr){4-5}\cmidrule(lr){6-8}\cmidrule(lr){9-10}
width & $\tau=1$ & annealed & $\Delta$ & $p$ & $\tau=1$ & annealed & $p$ & $\Delta$ & verdict \\
\midrule
%%MBTEMP%%
\bottomrule
\end{tabular}
\end{table}

\subsection{The pre-registered predictions}\label{sec:mbprereg}

Table~\ref{tab:mbprereg} gives the three predictions recorded before the
first MobileNetV2 run. P1 failed: the bank term was equivalent to zero in
\V{mbP1} of five conditions. The three inconclusive conditions are not
evidence of a bank effect, since none is significant, but the prediction
was equivalence and we score it as stated. P2 failed in \V{mbP2} of three
pairs, a reversal of the ResNet-18 result. P3 held in \V{mbP3} of five
conditions.

\begin{table}[!htbp]\centering\small
\caption{Predictions recorded on 30 September 2026, before any MobileNetV2
run, and the outcomes.}
\label{tab:mbprereg}
\begin{tabular}{l p{8.2cm} l l}
\toprule
& prediction & measured & outcome \\
\midrule
%%MBPREREG%%
\bottomrule
\end{tabular}
\end{table}

\section{What transfers between backbones}\label{sec:transfer}

Table~\ref{tab:contra} sets the main ResNet-18 results against
MobileNetV2.

\begin{table}[!htbp]\centering\small
\caption{ResNet-18 against MobileNetV2, result by result. Differences are
in pp and paired by seed.}
\label{tab:contra}
\resizebox{\textwidth}{!}{%
\begin{tabular}{l l l l}
\toprule
result & ResNet-18 & MobileNetV2 & status \\
\midrule
%%CONTRA%%
\bottomrule
\end{tabular}}
\end{table}

\paragraph{What holds.} On both backbones the bank adds nothing
measurable on its own, the routing term carries the effect wherever there
is one, and no bank size beyond $K=2$ is distinguishable from $K=2$. On
MobileNetV2 static and bank kernels share one initialisation, so its bank
term is free of the initialisation mismatch that affects the ResNet-18
bank term (Section~\ref{sec:threats}), and it gives the same answer.

\paragraph{Which cost to hold fixed.} The ResNet-18 conclusion that
parameters are better spent on width is a statement at equal parameters.
Dynamic convolution is usually advertised at equal compute. The two views
disagree (Table~\ref{tab:cost}). Doubling the width costs three to four
times the multiply-accumulates of the routed network; routing at the
same width costs at most \V{rbmacextra}\,\% more on ResNet-18 and
\V{mbmacextra}\,\% more on MobileNetV2, counted for an implementation that
aggregates the kernel per sample. At the same width, routing beats the
plain network in \V{rbeqmac} of five conditions on ResNet-18 and
\V{mbeqmac} of five on MobileNetV2. At equal parameters it loses all three
pairs on ResNet-18 and only one of three on MobileNetV2. Whether dynamic
convolution is worth its cost therefore depends on the cost and on the
backbone. Our training implementation computes $K$ convolutions per
dynamic layer and is slower in wall-clock time than either count
suggests (Tables~\ref{tab:config} and~\ref{tab:mbconfig}); that is a
property of the implementation, not of the layer.

\begin{table}[!htbp]\centering\small
\caption{The routed network under two cost views. Left: \texttt{static}
at width $2c$ minus \texttt{gap} at width $c$, with the parameter ratio
routed/plain and the multiply-accumulate (MAC) ratio plain/routed.
Right: \texttt{gap} minus \texttt{static} at the same width, with the MAC
ratio routed/plain. MACs per image are counted for per-sample kernel
aggregation.}
\label{tab:cost}
\resizebox{\textwidth}{!}{%
\begin{tabular}{l l r r l r l}
\toprule
& & \multicolumn{3}{c}{plain at $2c$ vs routed at $c$} & \multicolumn{2}{c}{routed vs plain, same width $c$} \\
\cmidrule(lr){3-5}\cmidrule(lr){6-7}
backbone & condition & param ratio & MAC ratio & plain $-$ routed (pp) & MAC ratio & routed $-$ plain (pp) \\
\midrule
%%COST%%
\bottomrule
\end{tabular}}
\end{table}

\paragraph{The CIFAR-100 width-64 null.} It does not appear on
MobileNetV2. Across the ten backbone and condition pairs, the only null
is ResNet-18 on CIFAR-100 at width 64, which is also the condition with
the lowest static training loss (\V{rblossCC}); no MobileNetV2 static arm
trains below \V{mblossCC} under this recipe (Figure~\ref{fig:losseffect}).
That pattern is what the saturation account of Section~\ref{sec:tin}
would expect, but ResNet-18 on Tiny-ImageNet at width 64 trains to a loss
of \V{lossTB} and shows one of the largest effects, so no single loss
threshold fits all ten points, and one chosen after seeing them would not
be evidence. We do not adopt it.

\begin{figure}[!htbp]\centering
\includegraphics[width=\textwidth]{fig/g6_loss_effect.pdf}
\caption{(a) MobileNetV2 training loss, seed 0, log scale. (b)
$\texttt{gap}-\texttt{static}$ against the static arm's final training
loss for both backbones and all five conditions, \SI{95}{\percent} CIs.}
\label{fig:losseffect}
\end{figure}

\paragraph{Temperature and the untrained router.} Annealing costs about
\SI{1}{pp} at every MobileNetV2 width (\V{mbanlist}\,pp at widths 16, 32
and 64), against \V{rbanlist}\,pp on ResNet-18. On MobileNetV2 the router
also does not need training. Both results point the same way: what
matters is that the mixture differ across images from the start of
training, not that it be optimised, and a schedule that keeps it flat
early removes most of the benefit. We state this as an interpretation;
neither experiment was designed to test it.

\paragraph{Why the backbones differ.} Several properties change at once
between ResNet-18 and MobileNetV2: the share of parameters in dynamic
layers (nearly all against \V{mbdynshare}\,\%), full against depthwise
convolution in those layers, the size of the router relative to the bank,
and how close the network comes to fitting its training set under this
recipe. None was isolated here. The three remaining backbones vary these
properties: VGG-16 keeps every parameter in full dynamic convolutions,
ResNet-50 about half, and ConvNeXt-Tiny makes only its depthwise
$7\times7$ layers dynamic.

\section{Threats to validity}\label{sec:threats}

\paragraph{Scale.} Both datasets are low-resolution and every run fits on
one consumer GPU. Behaviour at ImageNet scale is untested and cannot be
tested on this hardware.

\paragraph{Two backbones.} Two backbones disagree on three of the main
results, so what holds on both is the minimum that can be claimed. The
same program on VGG-16, ResNet-50 and ConvNeXt-Tiny was pre-registered on
30 September 2026 and is running; its results will be added as they
complete, whatever they show.

\paragraph{Placement and router size on MobileNetV2.} Copying the
ResNet-18 rule makes only the depthwise layers of MobileNetV2 dynamic,
\V{mbdynshare}\,\% of its parameters, and its router is larger than the
bank. A different rule, such as making every convolution but the stem
dynamic, could change the MobileNetV2 results. The frozen-router control
shows that the trainable router parameters do not produce the gain.

\paragraph{MobileNetV2 from the log.} The MobileNetV2 numbers are taken
from the training log's final lines, parsed twice by independent code
with identical results. The result files on the training machine record
the same values plus routing entropies; the entropy analyses
(Sections~\ref{sec:w64}, \ref{sec:temp} and~\ref{sec:entropy}) will be
repeated for MobileNetV2 from those files. One MobileNetV2 run was
resumed after a GPU fault; the conclusions do not depend on it.

\paragraph{One recipe.} Both backbones use the ResNet-18 recipe. The
MobileNetV2 static arms never approach zero training loss under it, and
a recipe tuned for MobileNetV2 could change where routing helps.

\paragraph{Initialisation.} Static convolutions in the \texttt{static}
arm, and the static stem, shortcut and head layers in every arm, keep
PyTorch's default initialisation (Kaiming-uniform with $a=\sqrt5$), while
every kernel of every dynamic bank is initialised Kaiming-normal with
fan-out. The \texttt{uniform}$-$\texttt{static} bank term therefore
compares two initialisations as well as a bank against a single kernel.
The routing term, \texttt{gap}$-$\texttt{uniform}, is not affected: both
arms have identically initialised banks. The replication on four more
backbones initialises static and bank kernels with one scheme, and the
MobileNetV2 bank term, which is free of the mismatch, also shows no bank
effect.

\paragraph{Not a reproduction of a published system.} The layer is a
dynamic convolution of the CondConv and DY-Conv form. No published
implementation of ODConv, FDConv or KernelWarehouse was run, and the
conclusions concern dynamic convolution of this form.

\paragraph{Global routing.} Every router here uses a global descriptor.
Spatially resolved routing is not covered.

\paragraph{One annealing schedule.} Section~\ref{sec:temp} compares a
constant temperature with one linear schedule.

\paragraph{Stored and effective parameters.} By
Proposition~\ref{prop:collapse} the \texttt{uniform} arm has the
effective capacity of a plain network. The matched-budget comparisons of
Section~\ref{sec:matched} use stored parameters. Section~\ref{sec:transfer}
gives the comparison at equal compute, which favours routing on both
backbones.

\paragraph{The CIFAR-100 width-64 null is unexplained.} It is
well-powered and does not depend on the temperature schedule, its stated
explanation failed a recorded test, and it does not appear on
MobileNetV2.

\paragraph{Provenance and corrections.} Final accuracies were transcribed
from the training logs, which print them verbatim, and every figure and
table is generated from those per-seed values. During the analysis, three
scripts grouped results by width alone, which merged runs from the two
datasets and from different bank sizes; this was detected when a
comparison paired one dataset's \texttt{uniform} arm with the other's
\texttt{static} arm, and every number in this report was recomputed from
per-seed values keyed on dataset, width, bank size and schedule. Summary
sentences printed by the analysis scripts were also found to over-read
non-significant orderings on four occasions; this report relies on the
per-comparison statistics only. The report generator was later
rewritten to read the result files directly; on the 24 width-64 result
files checked it reproduces the transcribed accuracies and routing
entropies exactly. The earlier version of this report printed stored
parameter counts with a third decimal that had not been measured, and six
of nineteen were off by up to \SI{0.005}{M} (for example \SI{0.708}{M}
for \SI{0.713}{M}); the counts here are counted from the model
definition and agree with the training logs to the two decimals they
print. No comparison changes.

\section{Conclusion}

On two backbones, the kernel bank of dynamic convolution adds nothing
measurable on its own. A uniformly mixed bank was never significantly
better or worse than a plain network, and was equivalent to one in
\V{bothbankeq} of ten conditions. Whatever dynamic convolution gains, it
gains through input-dependent mixing, and banks larger than $K=2$ buy
nothing further.

Whether that gain is worth having depends on the backbone and on the cost
held fixed. On ResNet-18, at equal parameters, the same parameters spent
on width do better by 2.3 to \SI{4.3}{pp}. On MobileNetV2 the routed
network is not distinguishable from a plain one with about twice its
parameters in two of three pairs. At equal compute, routing helps on
both. On MobileNetV2 the mixing does not have to be learned: a router
frozen at random initialisation does as well as a trained one, provided
each image is routed by its own descriptor.

Several results did not survive the second backbone. The ResNet-18 null
on CIFAR-100 at width 64 does not appear on MobileNetV2. Temperature
annealing, neutral or harmful on ResNet-18, costs about \SI{1}{pp} at
every MobileNetV2 width. Two of the three MobileNetV2 predictions failed,
as did the recorded ResNet-18 prediction for Tiny-ImageNet. The entropy
diagnostic, tested on ResNet-18, does not measure whether routing is
useful.

\paragraph{Reproducibility.} \V{totruns} runs (\V{nruns} ResNet-18,
\V{mbruns} MobileNetV2), \V{totgpuh} GPU-hours, one RTX 4070 Ti. Code, the pre-registration, the dated prediction, the gate
ledger with recorded evidence, and every per-seed value accompany this
report.

\FloatBarrier
\begin{thebibliography}{99}\small
\bibitem[Arora et~al.(2018)]{arora2018optimization} S.~Arora, N.~Cohen, and E.~Hazan. On the optimization of deep networks: Implicit acceleration by overparameterization. In \emph{ICML}, 2018.
\bibitem[Chen et~al.(2021)]{chen2021drconv} J.~Chen, X.~Wang, Z.~Guo, X.~Zhang, and J.~Sun. Dynamic region-aware convolution. In \emph{CVPR}, 2021.
\bibitem[Chen et~al.(2025)]{chen2025fdconv} L.~Chen, L.~Gu, L.~Li, C.~Yan, and Y.~Fu. Frequency dynamic convolution for dense image prediction. In \emph{CVPR}, pages 30178--30188, 2025.
\bibitem[Chen et~al.(2020)]{chen2020dynamic} Y.~Chen, X.~Dai, M.~Liu, D.~Chen, L.~Yuan, and Z.~Liu. Dynamic convolution: Attention over convolution kernels. In \emph{CVPR}, pages 11030--11039, 2020.
\bibitem[Ding et~al.(2019)]{ding2019acnet} X.~Ding, Y.~Guo, G.~Ding, and J.~Han. ACNet: Strengthening the kernel skeletons for powerful CNN via asymmetric convolution blocks. In \emph{ICCV}, 2019.
\bibitem[Ding et~al.(2021)]{ding2021repvgg} X.~Ding, X.~Zhang, N.~Ma, J.~Han, G.~Ding, and J.~Sun. RepVGG: Making VGG-style ConvNets great again. In \emph{CVPR}, 2021.
\bibitem[He et~al.(2023)]{he2023sdconv} S.~He, C.~Jiang, D.~Dong, and L.~Ding. SD-Conv: Towards the parameter-efficiency of dynamic convolution. In \emph{WACV}, 2023.
\bibitem[He et~al.(2016)]{he2016deep} K.~He, X.~Zhang, S.~Ren, and J.~Sun. Deep residual learning for image recognition. In \emph{CVPR}, 2016.
\bibitem[Krizhevsky(2009)]{krizhevsky2009} A.~Krizhevsky. Learning multiple layers of features from tiny images. Technical report, University of Toronto, 2009.
\bibitem[Lakens(2017)]{lakens2017} D.~Lakens. Equivalence tests: A practical primer for $t$ tests, correlations, and meta-analyses. \emph{Social Psychological and Personality Science}, 8(4):355--362, 2017.
\bibitem[Le and Yang(2015)]{le2015tiny} Y.~Le and X.~Yang. Tiny ImageNet visual recognition challenge. CS231N course report, Stanford University, 2015.
\bibitem[Li and Yao(2024)]{li2024kernelwarehouse} C.~Li and A.~Yao. KernelWarehouse: Rethinking the design of dynamic convolution. In \emph{ICML}, PMLR 235:29201--29221, 2024.
\bibitem[Li et~al.(2022)]{li2022odconv} C.~Li, A.~Zhou, and A.~Yao. Omni-dimensional dynamic convolution. In \emph{ICLR}, 2022.
\bibitem[Li et~al.(2021)]{li2021dcd} Y.~Li, Y.~Chen, X.~Dai, M.~Liu, D.~Chen, Y.~Yu, L.~Yuan, Z.~Liu, M.~Chen, and N.~Vasconcelos. Revisiting dynamic convolution via matrix decomposition. In \emph{ICLR}, 2021.
\bibitem[Sandler et~al.(2018)]{sandler2018mobilenetv2} M.~Sandler, A.~Howard, M.~Zhu, A.~Zhmoginov, and L.-C. Chen. MobileNetV2: Inverted residuals and linear bottlenecks. In \emph{CVPR}, 2018.
\bibitem[Schuirmann(1987)]{schuirmann1987} D.~J. Schuirmann. A comparison of the two one-sided tests procedure and the power approach for assessing the equivalence of average bioavailability. \emph{Journal of Pharmacokinetics and Biopharmaceutics}, 15(6):657--680, 1987.
\bibitem[Wilcoxon(1945)]{wilcoxon1945} F.~Wilcoxon. Individual comparisons by ranking methods. \emph{Biometrics Bulletin}, 1(6):80--83, 1945.
\bibitem[Yang et~al.(2019)]{yang2019condconv} B.~Yang, G.~Bender, Q.~V. Le, and J.~Ngiam. CondConv: Conditionally parameterized convolutions for efficient inference. In \emph{NeurIPS}, 2019.
\bibitem[Zhang et~al.(2017)]{zhang2017} C.~Zhang, S.~Bengio, M.~Hardt, B.~Recht, and O.~Vinyals. Understanding deep learning requires rethinking generalization. In \emph{ICLR}, 2017.
\end{thebibliography}

\clearpage
\appendix
\section{Complete results: ResNet-18}

\begin{figure}[h]\centering
\includegraphics[width=\textwidth]{fig/f2_allruns.pdf}
\caption{Every CIFAR-100 run at $K=4$ and $\tau=1$; bars are arm means.}
\label{fig:allruns}
\end{figure}

\begin{table}[h]\centering\scriptsize
\caption{Every run at $K=4$ and $\tau=1$: final-epoch top-1 accuracy
(\si{\percent}) for each seed, with mean and sample standard deviation.}
\label{tab:allruns}
\begin{tabular}{l r l rrrrrr rr}
\toprule
dataset & width & arm & s0 & s1 & s2 & s3 & s4 & s5 & mean & sd \\
\midrule
%%ALLRUNS%%
\bottomrule
\end{tabular}
\end{table}

\begin{table}[h]\centering\small
\caption{Paired comparisons against \texttt{static}, $n=6$. $p_W$ is the
Wilcoxon signed-rank test, which cannot fall below 0.031 at this $n$.}
\label{tab:vsstatic}
\begin{tabular}{l r l r l ll r l l}
\toprule
dataset & width & arm & $\Delta$ (pp) & \SI{95}{\percent} CI & $p_t$ & $p_W$ & $d_z$ & $p_{\text{TOST}}$ & verdict \\
\midrule
%%VSSTATIC%%
\bottomrule
\end{tabular}
\end{table}

\begin{table}[h]\centering\small
\caption{Configuration and cost. Router parameters are in thousands and
are included in the stored total. Training loss is the final-epoch mean
over six seeds.}
\label{tab:config}
\begin{tabular}{l r l rr r r}
\toprule
dataset & width & arm & params (M) & router (K) & train loss & min/run \\
\midrule
%%CONFIG%%
\bottomrule
\end{tabular}
\end{table}

\begin{table}[h]\centering\small
\caption{Routing entropy by arm, CIFAR-100 width 64, per-seed means.}
\label{tab:entw64}
\begin{tabular}{l rrr r}
\toprule
arm & mean $\tilde H$ & min & max & $1-\tilde H$ \\
\midrule
%%ENTW64%%
\bottomrule
\end{tabular}
\end{table}

\begin{table}[h]\centering\small
\caption{Every run in the study. GPU-hours are from logged wall-clock
time.}
\label{tab:invent}
\begin{tabular}{l r r l l r r}
\toprule
dataset & width & $K$ & $\tau$ & arms & runs & GPU-h \\
\midrule
%%INVENT%%
\bottomrule
\end{tabular}
\end{table}



\clearpage
\section{Complete results: MobileNetV2}

\begin{figure}[h]\centering
\includegraphics[width=\textwidth]{fig/g7_mb_allruns.pdf}
\caption{Every MobileNetV2 run at $K=4$ and $\tau=1$; bars are arm means.}
\label{fig:mballruns}
\end{figure}

\begin{table}[h]\centering\scriptsize
\caption{MobileNetV2, every run at $K=4$ and $\tau=1$: final-epoch top-1
accuracy (\si{\percent}). $^\dagger$ marks the run resumed after a GPU
fault.}
\label{tab:mballruns}
\begin{tabular}{l r l rrrrrr rr}
\toprule
dataset & width & arm & s0 & s1 & s2 & s3 & s4 & s5 & mean & sd \\
\midrule
%%MBALLRUNS%%
\bottomrule
\end{tabular}
\end{table}

\begin{table}[h]\centering\small
\caption{MobileNetV2: paired comparisons against \texttt{static}, $n=6$.}
\label{tab:mbvsstatic}
\resizebox{\textwidth}{!}{%
\begin{tabular}{l r l r l ll r l l}
\toprule
dataset & width & arm & $\Delta$ (pp) & \SI{95}{\percent} CI & $p_t$ & $p_W$ & $d_z$ & $p_{\text{TOST}}$ & verdict \\
\midrule
%%MBVSSTATIC%%
\bottomrule
\end{tabular}}
\end{table}

\begin{table}[h]\centering\small
\caption{MobileNetV2: configuration and cost. Parameters are counted from
the model definition; router parameters are in thousands (the frozen
router is held as buffers and not counted). Training loss is the
final-epoch mean over six seeds; minutes are logged wall-clock time.}
\label{tab:mbconfig}
\begin{tabular}{l r l rr r r}
\toprule
dataset & width & arm & params (M) & router (K) & train loss & min/run \\
\midrule
%%MBCONFIG%%
\bottomrule
\end{tabular}
\end{table}

\begin{table}[h]\centering\small
\caption{Every MobileNetV2 run. GPU-hours are from logged wall-clock time.}
\label{tab:mbinvent}
\begin{tabular}{l r r l l r r}
\toprule
dataset & width & $K$ & $\tau$ & arms & runs & GPU-h \\
\midrule
%%MBINVENT%%
\bottomrule
\end{tabular}
\end{table}

\end{document}
